% ---
% id: receiver/rem:exacteig
% kind: rem
% title: "The eigenvalue-slack assumption behind the exact window"
% summary: >-
%   In the exact-window maximum the analytic entry is a proven lower bound,
%   while the computed eigenvalue less the slack is one only under [[A-EIG]];
%   the rigorous evaluation replaces it by a proven bound. No test catches an
%   overshoot, r_6 detects only an undershoot, and each weight stays in [0, 1].
% prominence: standard
% status: accepted
% uses: [def:pmcex, thm:multiclick-bound, prop:exactnoworse, lem:exactbasic, lem:exactwindow, def:bounddata, lem:windowtail]
% mentions: [def:evaluation-model, def:slot, prop:slotvalid, receiver/detector-imperfections/rem:d2slack, def:tolmap, def:residuals, thm:emitted-bound, thm:exact-and-evaluated, req:exact-window-record, def:enclosed-weight, thm:enclosed-weight, req:rigorous-weight]
% assumes: [A-EIG]
% refs: []
% provenance: newly-derived
% lean:
% ---

The exact window trades the analytic bound for a sharper value on finitely
many blocks, and in the floating-point evaluation the price is one assumption
on the computed eigenvalues. The rigorous evaluation pays no such price: it
puts a rigorous lower bound on the block eigenvalue in the computed value's
place (Definition~\ref{def:enclosed-weight},
Requirement~\ref{req:rigorous-weight}), and the weight is then a bound by
Theorem~\ref{thm:enclosed-weight}.

\begin{remark}
\label{rem:exacteig}
\eqref{eq:elldef} is a maximum of two lower bounds on $\Bnd^{\ast(t)}(n)$,
and they are of two different \emph{kinds}. The first,
$\Bnd^{(t)}(n)\le\Bnd^{\ast(t)}(n)$, is Theorem~\ref{thm:multiclick-bound}: an
inequality between exact quantities, holding with no hypothesis about any
computation. The second,
$\widehat\Bnd^{\ast(t)}(n)-\varepsilon_{\mathrm{eig}}\le\Bnd^{\ast(t)}(n)$,
is \emph{not a theorem of this contract}:
the \tref{evaluation model} of Definition~\ref{def:evaluation-model} asserts no bound on
the difference between a quantity computed in floating point and the exact
one, and $\widehat\Bnd^{\ast}$ is the output
of an extreme-eigenvalue computation on a matrix of dimension
$\binom{|\InpC|+n-1}{n}$. It is Assumption~\ref{ass:eig}, and that
assumption is invoked \emph{only} at the settings and indices where a depth
is declared and the weight is that of Definition~\ref{def:pmcex}. What is a
theorem is the same maximum with a \tref[rigorous window bounds]{rigorous
window bound} in the second entry: in the \tref{enclosed weight} of
Definition~\ref{def:enclosed-weight} the entry is a real proved to lie below
$\Bnd^{\ast(t)}(n)$, not an evaluation of it, no slack is subtracted, and
Theorem~\ref{thm:enclosed-weight} makes the weight a bound with no assumption.
Everything below concerns the floating-point evaluation.

\emph{The direction.} A computed eigenvalue landing \emph{below} the exact
one costs tightness and nothing else, and \eqref{eq:elldef} floors the value
at $\Bnd^{(t)}(n)$, which is sound with no assumption at all; so that
direction cannot take the emitted weight below what the analytic evaluation
emits (Proposition~\ref{prop:exactnoworse}). A computed eigenvalue landing
\emph{above} the exact one by more than $\varepsilon_{\mathrm{eig}}$ can make
the emitted weight $\pmc^{(t)}$ of that setting exceed
$\inf_{n>N}\lmin(\blk{M}{\eqn}_{\mc|t})$: an over-large lower bound, the
unsafe direction, and one that reaches a key length through
Definition~\ref{def:slot}. No range check catches it, whatever the size of
the overshoot. Every branch of \eqref{eq:pmcex} is capped by an analytic
remainder that is at most $1$ and does not depend on
$\widehat\Bnd^{\ast(t)}$, so $\pmc^{(t)}$ stays in $[0,1]$ (the
\emph{Range} paragraph of this remark), and neither the ``$p_y<0$'' nor the
``$\sum_yp_y>1$'' failure condition that Proposition~\ref{prop:slotvalid}
names can fire on it. An overshoot of any size is caught by nothing. That
asymmetry is why Assumption~\ref{ass:eig} is
stated one-sidedly, on the side that matters, and why
$\varepsilon_{\mathrm{eig}}$ is \emph{subtracted} rather than used as a
comparison threshold: it decides nothing, it lowers a number.

\emph{A slack, not a repair.} The computed $\widehat\Bnd^{\ast(t)}(n)$
enters \eqref{eq:elldef} unmodified, and $\varepsilon_{\mathrm{eig}}$ enters
only as a subtracted declared slack, so that \eqref{eq:elldef} can be
recomputed from the recorded values
(Requirement~\ref{req:exact-window-record}). This is the distinction
Remark~\ref{rem:d2slack} draws for $\bar r$ and
Definition~\ref{def:tolmap}\,(iv) draws for $c$: a coefficient or a slack
inside a stated inequality is not a repair of a datum.

\emph{What $r_6$ tests, and the direction it is blind to.} Since
$\Bnd(n)\le\Bnd^{\ast}(n)$ is a theorem (Lemma~\ref{lem:exactbasic}), the
window consistency defect $r_6$ of \eqref{eq:res6} has exact value $0$, and
it is a residual in exactly the sense of the other five: what is emitted is
its floating-point evaluation, and a value above $\varepsilon_6$ refuses. It
catches an evaluation that contradicts an inequality this contract proves
--- a block formed at the wrong index, a wrongly ordered occupation basis, a
conditioned family taken at the wrong setting, an iteration that has not
converged --- which is the failure mode most likely to be gross. It does
\emph{not} verify Assumption~\ref{ass:eig}: $r_6$ fires when
$\widehat\Bnd^{\ast}$ lands too \emph{low}, and Assumption~\ref{ass:eig}
constrains it from landing too \emph{high}. No test in this contract bounds
the unsafe direction; the assumption carries it.

\emph{Range.} $0\le\pmc^{(t)}\le1$ holds exactly, whatever the computed
values $\widehat\Bnd^{\ast(t)}(n)$ and whether or not
Assumption~\ref{ass:eig} holds: the outer maximum of \eqref{eq:pmcex} gives
$\pmc^{(t)}\ge0$, and the inner minimum of each branch contains an analytic
remainder at most $1$ --- $\Bnd^{(t)}(m^{\ast}_t+1)\le1$ on branch~1 by
Theorem~\ref{thm:multiclick-bound}, and the tail value
$T^{(t)}_{m_t}=1-\sum_k\beta^{(t)}_k\big(d^{(t)}_k\big)^{m_t}\le1$ of
\eqref{eq:windowtail} on branch~2 because $\beta^{(t)}_k,d^{(t)}_k\ge0$ by
\eqref{eq:bounddata}. Lemma~\ref{lem:exactwindow} gives the same range only
through \eqref{eq:ellhyp}, which at a setting carrying a depth is
Assumption~\ref{ass:eig}. So Theorem~\ref{thm:emitted-bound}'s range statement
and Proposition~\ref{prop:slotvalid} hold unchanged; their floating-point
evaluation is governed by Theorem~\ref{thm:exact-and-evaluated}\,(ii) here as everywhere
else, and no computed value is clamped into $[0,1]$ to make them look true.
\end{remark}
